%%%%%%%%%%%%%%%%%%%%%%%%%%%%%%%%%%%%%%%%%%%%%%%%%%%%%%
%%%%               "Front matter":               %%%%%
%%%%                                             %%%%%
%%%%  - Cover page (Titelblatt)                  %%%%%
%%%%  - Curriculum Vitae                         %%%%%
%%%%  - Abstract (English)                       %%%%%
%%%%  - Kurzzusammenfassung (German)             %%%%%
%%%%  - Acknowledgements                         %%%%%
%%%%  - Table of contents                        %%%%%
%%%%  - Notation / Symbols                       %%%%%
%%%%%%%%%%%%%%%%%%%%%%%%%%%%%%%%%%%%%%%%%%%%%%%%%%%%%%
%
%%%%%%%%%%%%%%%%%%%%%%%%%%%%%%%%%%%%%%%%%%%%%%%%%%%%%%
%% "Magic comments" for the LaTeX editor -------------
% !TEX root = ../Diss.tex
%%%%%%%%%%%%%%%%%%%%%%%%%%%%%%%%%%%%%%%%%%%%%%%%%%%%%%
%
%%%%%%%%%%%%%%%%%%%%%%%%%%%%%%%%%%%%%%%%
%%%%           Deckblatt            %%%%
%%%%%%%%%%%%%%%%%%%%%%%%%%%%%%%%%%%%%%%%
\begin{showCoverPage}%
%% define an (invisible) numbering for the cover page
%% that is different from those used in the document
%% (i.e. Roman and Arabic) in order to prevent warnings
%% due to duplicates in page numbering
%% ("destinatiion with same identifier (name{page.1})...")
\pagenumbering{alph}%
%% Set a PDF bookmark for the cover page:
\pdfbookmark[0]{\TransTitlePageName}{cover}%
%% Cover page:
\begin{titlepage}
\thispagestyle{empty}

%%%%Dazu muss \usepackage[absolute]{textpos} eingebunden werden.
%%%%Die erste Zahl gibt die Breite des Textblocks an (wird fur Bilder NICHT!!! ignoriert, legt die Bildbreite fest), die
%%%%zweite den Abstand zum linken und die dritte den Abstand zum oberen Seitenrand.
%
\mbox{}%
%%%%%%%%%%%%  TITEL DER ARBEIT
 	\begin{textblock}{11}[0.5,0.5](7.8,4.5)
		\sffamily%
		%\huge%
		\LARGE%
 		\centering%
		\textbf{\ThesisTitleDivided}
 	\end{textblock}

%%%%%%%%%%%% TYP DER ARBEIT
%%%%%%%%%%%% AUTORNAME UND DATUM
 	\begin{textblock}{11}[0.5,0.5](7.8,11.0)
 		\sffamily%
		\centering%
		zur Erlangung des akademischen Grades eines\\%
		\vspace{0.5\baselineskip}%
		\Large%
		\DocDegree%
		\vspace{1\baselineskip}\\%
		\normalsize%
		\ifthenelse{\boolean{statusApproved}}{von der}{der}
		\Department\\%
		\UniversityGenitiv%
		\vspace{0.75\baselineskip}\\%
		\textbf{\ifthenelse{\boolean{statusApproved}}{genehmigte}{vorgelegte}}%
		\vspace{0.5\baselineskip}\\%
		{\Large\TypeOfThesis}%
		\vspace{0.5\baselineskip}\\%
		\normalsize%
		von%
		\vspace{0.5\baselineskip}\\%
		\begin{showCurrentAutorTitle}%
		%\textbf{\Large%               <----- Einkommentieren, falls der Titel mit dem Namen in einer Zeile angezeigt werden soll
		\AuthorTitle
		%}%                            <----- Einkommentieren, falls der Titel mit dem Namen in einer Zeile angezeigt werden soll
		\vspace{0.5\baselineskip}\\%  <----- Auskommentieren, falls der Titel mit dem Namen in einer Zeile angezeigt werden soll
		\end{showCurrentAutorTitle}%
		\textbf{\Large%
		\AuthorName}%
		\vspace{0.5\baselineskip}\\%
		\normalsize%
		\begin{showPlaceOfBirth}%
		aus \PlaceOfBirth%
		\end{showPlaceOfBirth}%
 	\end{textblock}

%%%%%%%%%%%% BETREUER
	\begin{textblock}{11}[0.5,0.5](7.8,17)
	%\begin{textblock}{9}(2.4,17)
	\sffamily%
	\raggedright%
	\setlength{\tabcolsep}{0pt}
		\begin{tabularx}{\columnwidth}{x@{\extracolsep{20pt}}l}
			Tag der mündlichen Prüfung: & \ExamDate\texorpdfstring{\\}{}
			Erster Gutachter: & \MainAdvisor\texorpdfstring{\\}{}
			Zweiter Gutachter: & \CoAdvisor\texorpdfstring{\\}{}
		\end{tabularx}
 	\end{textblock}
	
\end{titlepage}
\clearpage
\end{showCoverPage}%
\frontmatter
%
%
%%%%%%%%%%%%%%%%%%%%%%%%%%%%%%%%%%%%%%%%
%%%% Kurzeinführung in die Vorlage  %%%%
%%%%%%%%%%%%%%%%%%%%%%%%%%%%%%%%%%%%%%%%
\begin{showExamples}%
%
%%%%%%%%%%%%%%%%%%%%%%%%%%%%%%%%%%%%%%%%%%%%%%%%%%%%%%
%% "Magic comments" for the LaTeX editor -------------
% !TEX root = ../Diss.tex
%%%%%%%%%%%%%%%%%%%%%%%%%%%%%%%%%%%%%%%%%%%%%%%%%%%%%%
%
%%%%%%%%%%%%%%%%%%%%%%%%%%%%%%%%%%%%%
\chapter*{Kurzeinführung in diese Vorlage}
%%%%%%%%%%%%%%%%%%%%%%%%%%%%%%%%%%%%%
%
%% Add an entry for this chapter to the table of contents:
\addcontentsline{toc}{chapter}{Kurzeinführung in die Vorlage}
%
%% If no entry to the table of contents:
%% Set a PDF bookmark (PDF-Lesezeichen):
%\pdfbookmark[0]{Kurzeinführung in die Vorlage}{kurzeinfuehrung}
%
%Add a header title (Kolumnentitel beidseitig setzen)
\markboth{Kurzeinführung in die Vorlage}{Kurzeinführung in die Vorlage}
%
%%%%%%%%%%%%%%%%%%%%%%%%%%%%%%%%%%%%%
%
Herzlich Willkommen zur \LaTeX-Vorlage
%Nutzung von \glsdisp wegen der Genitiv-Form
des \glsgen{ac:IES} %{Lehrstuhls für Interalktive Echtzeitsysteme (IES)}
des \glsgen{ac:KIT}. %{Karlsruher Institus für Technilogie (KIT)}.
Sie wurde ursprünglich von
\href{https://ies.iar.kit.edu/1473_1575.php}{Philipp Woock}
entwickelt und basiert in ihren Grundzügen auf der \enquote{allgemeinen sehr umfassenden}
\href{http://www.matthiaspospiech.de/latex/templates/thesis/}{Vorlage} von 
\href{http://www.matthiaspospiech.de}{Matthias Pospiech} von der Leibniz Universität Hannover.
Ohne diese Basis wäre die Vorlage niemals das geworden, was sie ist. Vielen Dank!

Im Jahre 2018 wurde diese Vorlage von
\href{https://ies.iar.kit.edu/1473_1650.php}{Michael Grinberg}
grundlegend überarbeitet und an die neuen Vorgaben des \glsgen{ac:KSP} angepasst.
Insbesondere wurde die Vorlage dahingehend abgeändert, dass sie nun XeLaTeX statt pdfLaTeX, 
\gls{gls:biblatex} statt \gls{gls:bibtex} verwendet und die inzwischen aufgetretenen
Inkompatibilitäten mit neuen Versionen einiger \gls{gls:latex}-\glspl{gls:pkg} behebt.
In 2019 wurde die Vorlage weiter angepasst und mit Beispielen aus einer Vorlage
für studentische Arbeiten, die von
\href{https://crypto.iti.kit.edu/alumni_matthias_nagel.php}{Mathias Nagel}
entwickelt worden ist, angereichert.

Zur Nutzung dieser Vorlage für die eigene Arbeit empfiehlt es sich, 
die Anleitungskapitel auszublenden und ansonsten die vordefinierte Struktur zu nutzen.
Die Anleitungskapitel lassen sich bequem ausblenden, indem man in der Datei
\printfilepathquoted{preambel/AlleSchalter.tex} in der Zeile
\lstinline|\setUserDefinedBoolean{showExamples}{true}| 
den Wert \printkeyword{true} durch \printkeyword{false} ersetzt.
Zur Einblendung der vordefinierten Kapitelstruktur
sollte stattdessen der Wert der Variablen
\printkeyword{showIntro},
\printkeyword{showRelatedWork},
\printkeyword{showConcept}
\usw
%\printkeyword{showSystem},
%\printkeyword{showEvaluation},
%\printkeyword{showConclusion} und
%\printkeyword{showOutlook}
auf \printkeyword{true} gesetzt werden.

\begin{tcolorbox}[colback=orange]
\footnotesize
Hinweis:
Zu beachten ist, dass zum Kompilieren der Vorlage
\printkeyword{xelatex} statt \printkeyword{pdflatex}
und zum Kompilieren der Bibliographien
\printkeyword{biber} statt \printkeyword{bibtex}
verwendet werden sollen.
Zur Generierung des Begriffs- und Abkürzungsverzeichnisses muss zusätzlich
\printkeyword{makeglossaries} (benötigt Perl!) aufgerufen werden.
\end{tcolorbox}

%
\cleardoublepage
\end{showExamples}%
%
%
%%%%%%%%%%%%%%%%%%%%%%%%%%%%%%%%%%%%%%%%
%%%%    Lebenslauf (deutsch)        %%%%
%%%%%%%%%%%%%%%%%%%%%%%%%%%%%%%%%%%%%%%%
\begin{showGermanCV}%
%
%%%%%%%%%%%%%%%%%%%%%%%%%%%%%%%%%%%%%
\chapter*{Lebenslauf}
%%%%%%%%%%%%%%%%%%%%%%%%%%%%%%%%%%%%%
%
%% No entry to the table of contents
%\addcontentsline{toc}{chapter}{Curriculum Vitae}
%
%% However, set a PDF bookmark:
\pdfbookmark[0]{Lebenslauf}{cv-de}
%
%Add a header title (Kolumnentitel setzen)
\markboth{Lebenslauf}{Lebenslauf}
%
%%%%%%%%%%%%%%%%%%%%%%%%%%%%%%%%%%%%%
%
%
%%%%%%%%%%%%%%%%%%%%%%%%%%%%%%%%%%%%%%%%%%%%%%%%%%%%%%%%%%%%
\todo[inline,caption={Deutschsprachigen Lebenslauf einfügen oder diese Seite ausblenden!}]{%
Beim Vorlegen der Dissertation an der KIT-Fakultät für Informatik muss ein Lebenslauf beigefügt werden.
Wird diese Seite nicht benötigt (\zB wenn man eine eigene Vorlage verwendet
oder wenn die Dissertation veröffentlicht wird), lässt sich diese Seite
durch die Umstellung des Schalters \printkeyword{showGermanCV}
auf \printkeyword{false} in der Datei \printfilepathquoted{./preambel/AlleSchalter.tex}
komplett ausblenden.
}%
%%%%%%%%%%%%%%%%%%%%%%%%%%%%%%%%%%%%%%%%%%%%%%%%%%%%%%%%%%%%
%
%
%%%%%%%%%%%%%%%%%%%%%%%%%%%%%%%%%%%%%%%%%%%%%%%%%%%%%%%%%%%%
\minisec{Persönliche Angaben}
%%%%%%%%%%%%%%%%%%%%%%%%%%%%%%%%%%%%%%%%%%%%%%%%%%%%%%%%%%%%
%% Durch Definition einer Gruppe mit \begingroup ... \endgroup 
%% sicherstellen, dass Anpassungen nur lokal wirken
%% und sich nicht auf die nachfolgenden tabellen erstrecken
\begingroup%
\small% Kleinere Schrift
\setlength{\tabcolsep}{0pt}% Kein Einzug
\renewcommand{\arraystretch}{1.0}% Abstand zwischen den Zeilen nicht vergrößern
\begin{longtable}{>{\raggedright\arraybackslash}p{0.4\linewidth-2\tabcolsep}>{\raggedright\arraybackslash}p{0.6\linewidth-2\tabcolsep}}%
%%-----------------------------------------------------------------------------------------------------
Titel: & \AuthorTitle \\
Name: & \AuthorName \\
Geburtsdatum: & 01.01.1990 \\
Geburtsort: & \PlaceOfBirth \\
Nationalität: & deutsch \\
%%-----------------------------------------------------------------------------------------------------
\end{longtable}
\endgroup%

%%%%%%%%%%%%%%%%%%%%%%%%%%%%%%%%%%%%%%%%%%%%%%%%%%%%%%%%%%%%
\minisec{Bildung}
%%%%%%%%%%%%%%%%%%%%%%%%%%%%%%%%%%%%%%%%%%%%%%%%%%%%%%%%%%%%
%% Durch Definition einer Gruppe mit \begingroup ... \endgroup 
%% sicherstellen, dass Anpassungen nur lokal wirken
%% und sich nicht auf die nachfolgenden tabellen erstrecken
\begingroup%
\small% Kleinere Schrift
\setlength{\tabcolsep}{0pt}% Kein Einzug
\renewcommand{\arraystretch}{1.3}% Etwas mehr Abstand zwischen den Zeilen, damit sie nicht zusammenfallen
\begin{longtable}{>{\raggedright\arraybackslash}p{0.4\linewidth-2\tabcolsep}>{\raggedright\arraybackslash}p{0.6\linewidth-2\tabcolsep}}%
%%-----------------------------------------------------------------------------------------------------
01.01.1900 -- 01.01.2000 & \textbf{Gymnasium XYZ in ABC}\\
%%-----------------------------------------------------------------------------------------------------
01.01.1900 -- 01.01.2000 & \textbf{Studium der XXX am KIT}\newline Abschluss als ...\\
%%-----------------------------------------------------------------------------------------------------
\end{longtable}
\endgroup%

%%%%%%%%%%%%%%%%%%%%%%%%%%%%%%%%%%%%%%%%%%%%%%%%%%%%%%%%%%%%
\minisec{Berufserfahrung}
%%%%%%%%%%%%%%%%%%%%%%%%%%%%%%%%%%%%%%%%%%%%%%%%%%%%%%%%%%%%
%% Durch Definition einer Gruppe mit \begingroup ... \endgroup 
%% sicherstellen, dass Anpassungen nur lokal wirken
%% und sich nicht auf die nachfolgenden tabellen erstrecken
\begingroup%
\small% Kleinere Schrift
\setlength{\tabcolsep}{0pt}% Kein Einzug
\renewcommand{\arraystretch}{1.3}% Etwas mehr Abstand zwischen den Zeilen, damit sie nicht zusammenfallen
\begin{longtable}{>{\raggedright\arraybackslash}p{0.4\linewidth-2\tabcolsep}>{\raggedright\arraybackslash}p{0.6\linewidth-2\tabcolsep}}%
%-----------------------------------------------------------------------------------------------------
01.01.1900 -- 01.01.2000 & \textbf{Stelle XXX}\newline Wissenschaftliche Hilfskraft\\
%-----------------------------------------------------------------------------------------------------
01.01.1900 -- 01.01.2000 & \textbf{Stelle YYY}\newline Berufsbezeichnung\\
%%-----------------------------------------------------------------------------------------------------
\end{longtable}
\endgroup%

%
\clearpage
\end{showGermanCV}%
%
%
%%%%%%%%%%%%%%%%%%%%%%%%%%%%%%%%%%%%%%%%
%%%%  Curriculum Vitae (Englisch)   %%%%
%%%%%%%%%%%%%%%%%%%%%%%%%%%%%%%%%%%%%%%%
\begin{showEnglishCV}%
%
%% Set English as the text language
%% independent on the document language
\textInEnglish{%
%%%%%%%%%%%%%%%%%%%%%%%%%%%%%%%%%%%%%
\chapter*{Curriculum Vitae}
%%%%%%%%%%%%%%%%%%%%%%%%%%%%%%%%%%%%%
%
%% No entry to the table of contents
%\addcontentsline{toc}{chapter}{Curriculum Vitae}
%
%% However, set a PDF bookmark:
\pdfbookmark[0]{Curriculum Vitae}{cv-en}
%
%Add a header title (Kolumnentitel setzen)
\markboth{Curriculum Vitae}{Curriculum Vitae}
%
%%%%%%%%%%%%%%%%%%%%%%%%%%%%%%%%%%%%%
%
%
%%%%%%%%%%%%%%%%%%%%%%%%%%%%%%%%%%%%%%%%%%%%%%%%%%%%%%%%%%%%
\todo[inline,caption={Add your curriculum vitae in English or hide the corresponding page!}]{%
English version of the CV if needed.

This page can be hidden by setting the switch \printkeyword{showEnglishCV}
to \printkeyword{false} in the file \printfilepathquoted{./preambel/AlleSchalter.tex}.
}%
%%%%%%%%%%%%%%%%%%%%%%%%%%%%%%%%%%%%%%%%%%%%%%%%%%%%%%%%%%%%
%
%
%%%%%%%%%%%%%%%%%%%%%%%%%%%%%%%%%%%%%%%%%%%%%%%%%%%%%%%%%%%%
\minisec{Personal details}
%%%%%%%%%%%%%%%%%%%%%%%%%%%%%%%%%%%%%%%%%%%%%%%%%%%%%%%%%%%%
%% Durch Definition einer Gruppe mit \begingroup ... \endgroup 
%% sicherstellen, dass Anpassungen nur lokal wirken
%% und sich nicht auf die nachfolgenden tabellen erstrecken
\begingroup%
\small% Kleinere Schrift
\setlength{\tabcolsep}{0pt}% Kein Einzug
\renewcommand{\arraystretch}{1.0}% Abstand zwischen den Zeilen nicht vergrößern
\begin{longtable}{>{\raggedright\arraybackslash}p{0.4\linewidth-2\tabcolsep}>{\raggedright\arraybackslash}p{0.6\linewidth-2\tabcolsep}}%
%%-----------------------------------------------------------------------------------------------------
Title: & \AuthorTitle \\
Name: & \AuthorName \\
Date of birth: & Jan. 01, 1990 \\
Place of Birth: & \PlaceOfBirth \\
Nationality: & German \\
%%-----------------------------------------------------------------------------------------------------
\end{longtable}
\endgroup%

%%%%%%%%%%%%%%%%%%%%%%%%%%%%%%%%%%%%%%%%%%%%%%%%%%%%%%%%%%%%
\minisec{Education}
%%%%%%%%%%%%%%%%%%%%%%%%%%%%%%%%%%%%%%%%%%%%%%%%%%%%%%%%%%%%
%% Durch Definition einer Gruppe mit \begingroup ... \endgroup 
%% sicherstellen, dass Anpassungen nur lokal wirken
%% und sich nicht auf die nachfolgenden tabellen erstrecken
\begingroup%
\small% Kleinere Schrift
\setlength{\tabcolsep}{0pt}% Kein Einzug
\renewcommand{\arraystretch}{1.3}% Etwas mehr Abstand zwischen den Zeilen, damit sie nicht zusammenfallen
\begin{longtable}{>{\raggedright\arraybackslash}p{0.4\linewidth-2\tabcolsep}>{\raggedright\arraybackslash}p{0.6\linewidth-2\tabcolsep}}%
%%-----------------------------------------------------------------------------------------------------
Jan. 01, 1990 -- Jan. 01, 2000 & \textbf{Schulbezeichnung in ABC}\\
%%-----------------------------------------------------------------------------------------------------
Jan. 01, 1990 -- Jan. 01, 2000 & \textbf{Studiumsbezeichnung of XXX at YYY}\newline Abshluss\\
%%-----------------------------------------------------------------------------------------------------
\end{longtable}
\endgroup%

%%%%%%%%%%%%%%%%%%%%%%%%%%%%%%%%%%%%%%%%%%%%%%%%%%%%%%%%%%%%
\minisec{Professional experience}
%%%%%%%%%%%%%%%%%%%%%%%%%%%%%%%%%%%%%%%%%%%%%%%%%%%%%%%%%%%%
%% Durch Definition einer Gruppe mit \begingroup ... \endgroup 
%% sicherstellen, dass Anpassungen nur lokal wirken
%% und sich nicht auf die nachfolgenden tabellen erstrecken
\begingroup%
\small% Kleinere Schrift
\setlength{\tabcolsep}{0pt}% Kein Einzug
\renewcommand{\arraystretch}{1.3}% Etwas mehr Abstand zwischen den Zeilen, damit sie nicht zusammenfallen
\begin{longtable}{>{\raggedright\arraybackslash}p{0.4\linewidth-2\tabcolsep}>{\raggedright\arraybackslash}p{0.6\linewidth-2\tabcolsep}}%
%-----------------------------------------------------------------------------------------------------
Jan. 01, 1990 -- Jan. 01, 2000 & \textbf{Stelle XXX}\newline Berufsbezeichnung\\
%-----------------------------------------------------------------------------------------------------
Jan. 01, 1990 -- Jan. 01, 2000 & \textbf{Stelle YYY}\newline Berufsbezeichnung\\
%%-----------------------------------------------------------------------------------------------------
\end{longtable}
\endgroup%

}%endtextInEnglish
%
\clearpage
\end{showEnglishCV}%
%
%
%%%%%%%%%%%%%%%%%%%%%%%%%%%%%%%%%%%%%%%%
%%%%       Abstract (English)       %%%%
%%%%%%%%%%%%%%%%%%%%%%%%%%%%%%%%%%%%%%%%
\begin{showAbstractEN}%
%
%% Set English as the text language
%% independent on the document language
\textInEnglish{%
%%%%%%%%%%%%%%%%%%%%%%%%%%%%%%%%%%%%%
\chapter*{Abstract}
%%%%%%%%%%%%%%%%%%%%%%%%%%%%%%%%%%%%%
%
%% Add an entry to the table of contents and a PDF bookmark
\addcontentsline{toc}{chapter}{Abstract}
%
%% If no entry to the table of contents:
%% Set a PDF bookmark (PDF-Lesezeichen):
%\pdfbookmark[0]{Abstract}{abstract}
%
%Add a header title (Kolumnentitel setzen)
\markboth{Abstract}{Abstract}
%
%%%%%%%%%%%%%%%%%%%%%%%%%%%%%%%%%%%%%
%
%\todo[inline]{Add the English abstract of the thesis here.}
%
%%%%%%%%%%%%%%%%%%%%%%%%%%%%%%%%%%%%%
Add the English abstract of the thesis here.

\Blindtext
}%endtextInEnglish
%
\clearpage
%
\end{showAbstractEN}%
%
%
%%%%%%%%%%%%%%%%%%%%%%%%%%%%%%%%%%%%%%%%
%%%% Kurzzusammenfassung (Deutsch)  %%%%
%%%%%%%%%%%%%%%%%%%%%%%%%%%%%%%%%%%%%%%%
\begin{showAbstractDE}%
%
%% Set German as the text language
%% independently on the document language
\textInGerman{%
%%%%%%%%%%%%%%%%%%%%%%%%%%%%%%%%%%%%%
\chapter*{Kurzfassung}
%%%%%%%%%%%%%%%%%%%%%%%%%%%%%%%%%%%%%
%
%% Add an entry to the table of contents and a PDF-Bookmark
\addcontentsline{toc}{chapter}{Kurzfassung}
%
%% If no entry to the table of contents:
%% Set a PDF bookmark (PDF-Lesezeichen):
%\pdfbookmark[0]{Kurzfassung}{kurzzusammenfassung}
%
%Add a header title (Kolumnentitel setzen)
\markboth{Kurzfassung}{Kurzfassung}
%
%%%%%%%%%%%%%%%%%%%%%%%%%%%%%%%%%%%%%
%
%\todo[inline]{Hier die deutsche Kurzfassung einfügen.}
%
%%%%%%%%%%%%%%%%%%%%%%%%%%%%%%%%%%%%%
Hier die deutsche Kurzfassung einfügen.

\Blindtext
}%endtextInGerman
%
\cleardoublepage
\end{showAbstractDE}%
%
%
%%%%%%%%%%%%%%%%%%%%%%%%%%%%%%%%%%%%%%%%
%%%%           Danksagung           %%%%
%%%%%%%%%%%%%%%%%%%%%%%%%%%%%%%%%%%%%%%%
\begin{showAcknowledgments}%
%
%%%%%%%%%%%%%%%%%%%%%%%%%%%%%%%%%%%%%
% Note: the translation macros (e.g. "\TransAcknowledgements")
% are to be found in /preambel/Translations.tex
%%%%%%%%%%%%%%%%%%%%%%%%%%%%%%%%%%%%%
\chapter*{\TransAcknowledgements}
%%%%%%%%%%%%%%%%%%%%%%%%%%%%%%%%%%%%%
%
%% Add an entry to the table of contents and a PDF bookmark
\addcontentsline{toc}{chapter}{\TransAcknowledgements}
%
%% If no entry to the table of contents:
%% Set a PDF bookmark (PDF-Lesezeichen):
%\pdfbookmark[0]{\TransAcknowledgements}{acknowledgements}
%
%Add a header title (Kolumnentitel beidseitig setzen)
\markboth{\TransAcknowledgements}{\TransAcknowledgements}
%
%%%%%%%%%%%%%%%%%%%%%%%%%%%%%%%%%%%%%
%
%%%%%%%%%%%%%%%%%%%%%%%%%%%%%%%%%%%%%%%%%%%%%%%%%%%%%%%%%%%%
\todo[inline,caption={Danksagung einfügen oder ausblenden!}]{%
TODO: Danksagung einfügen!

Notiz: Die Seite kann durch das Umsetzen des Schalters
\printkeyword{showAcknowledgments} auf \printkeyword{false}
in der Datei \printfilepathquoted{preambel/AlleSchalter.tex}
ausgeblendet werden.
}%
%%%%%%%%%%%%%%%%%%%%%%%%%%%%%%%%%%%%%%%%%%%%%%%%%%%%%%%%%%%%
Hier steht die Danksagung.

\blindtext
%
\cleardoublepage
\end{showAcknowledgments}%
%
%
%%%%%%%%%%%%%%%%%%%%%%%%%%%%%%%%%%%%%%%%
%%%%       Inhaltsverzeichnis       %%%%
%%%%%%%%%%%%%%%%%%%%%%%%%%%%%%%%%%%%%%%%
\begin{showTableOfContents}%
%\cleardoublepage
%% Inhaltsverzeichnis in die PDF-Links eintragen
%% damit es als Lesezeichen im PDF auftaucht
\phantomsection
\pdfbookmark[0]{\protect\contentsname}{toc}
%% Um günstige Seitenumbrücke zu bewirken,
%% kann es sinnvoll sein, die Zeilenabstände
%% im Inhaltsverzeichnis zu vergrößern.
%% Dies geht mit dem Makro \linespread{faktor}
%% (Faktor kann angepasst werden):
%\linespread{1.25}
\hypertarget{toc}{}%make toc linkable (before tableofcontents!)
\tableofcontents
\clearpage
\end{showTableOfContents}%
%
%
%
%%%%%%%%%%%%%%%%%%%%%%%%%%%%%%%%%%%%%%%%
%%%%            Notation            %%%%
%%%%     bzw. manuell erstelltes    %%%%
%%%%       Symbolverzeichnis        %%%%
%%%%%%%%%%%%%%%%%%%%%%%%%%%%%%%%%%%%%%%%
\begin{showNotation}
%%%%%%%%%%%%%%%%%%%%%%%%%%%%%%%%%%%%%%%%%%%%%%%%%%%%%%%%%%%%
\chapter*{Notation}
%%%%%%%%%%%%%%%%%%%%%%%%%%%%%%%%%%%%%%%%%%%%%%%%%%%%%%%%%%%%
%
%% Add an entry to the table of contents and a PDF bookmark
\addcontentsline{toc}{chapter}{Notation}
%
%% No need for a separate bookmark (done via toc entry above)
%\pdfbookmark[0]{Notation}{notation}
%
%Add a header title (Kolumnentitel setzen)
\markboth{Notation}{Notation}
%
This chapter introduces the notation and symbols which are used in this thesis.
%In cases where a symbol has more than one meaning, the context (or a specific statement) resolves the ambiguity.

%%%%%%%%%%%%%%%%%%%%%%%%%%%%%%%%%%%%%%%%%%%%%%%%%%%%%%%%%%%%
\section*{General notation}
%%%%%%%%%%%%%%%%%%%%%%%%%%%%%%%%%%%%%%%%%%%%%%%%%%%%%%%%%%%%
\begin{myNotationDescTable}[0.3\linewidth] %Breite der ersten Spalte anpassbar, sollte aber überall gleich sein
%
% Als dreispaltige Tabelle angelegt,
% Beim Makro myNotationDescTableEntry ist die dritte Spalte bereits im Mathe-Modus
\myNotationDescTableEntry{Scalars}{italic Roman and Greek lowercase letters}{x, \alpha}
\myNotationDescTableEntry{Sets}{Greek uppercase letters}{\Theta}
\myNotationDescTableEntry{Vectors}{bold Roman lowercase letters}{\vt}
\myNotationDescTableEntry{Matrices}{bold Roman uppercase letters}{\mR}
\myNotationDescTableEntry{State spaces}{bold calligraphic Roman uppercase letters}{\spX}
\myNotationDescTableEntry{Random variables}{italic Roman uppercase letters}{\rvE}
\myNotationDescTableEntry{Multi-dimensional random variables}{bold italic Roman uppercase letters}{\mdrvE}
%\myNotationDescTableEntry{Distributions}{calligraphic uppercase letters }{\mathcal{N}}
%
\end{myNotationDescTable}
\vspace{-\parskip}
In multidimensional sets of elements related to time series, the first index denotes time.
%
%%%%%%%%%%%%%%%%%%%%%%%%%%%%%%%%%%%%%%%%%%%%%%%%%%%%%%%%%%%%
%%
%% Ggf. kann ein eigenes Symbolverzeichnis ("List of Symbols") am
%% Ende des Manuskripts eingefügt werden (s. Inhalt-BackMatter.tex).
%% Um dieses zu erzeugen, müssen Symbole in Form von Glossareinträgen
%% angelegt und im Text der Arbeit entsprechend referenziert werden.
%% Normalerweise reicht jedoch eine manuell erstellte Übersicht
%% über die einzelnen Symbole wie nachfolgend angeführt.
%%
%%%%%%%%%%%%%%%%%%%%%%%%%%%%%%%%%%%%%%%%%%%%%%%%%%%%%%%%%%%%

%%%%%%%%%%%%%%%%%%%%%%%%%%%%%%%%%%%%%%%%%%%%%%%%%%%%%%%%%%%%
\section*{Distributions}
%%%%%%%%%%%%%%%%%%%%%%%%%%%%%%%%%%%%%%%%%%%%%%%%%%%%%%%%%%%%
\begin{myNotationTable}[0.3\linewidth] %Breite der ersten Spalte anpassbar, sollte aber überall gleich sein
%
% Als zweispaltige Tabelle angelegt,
% Beim Makro myNotationTableEntryMath ist die erste Spalte bereits im Mathe-Modus
\myNotationTableEntryMath{\GaussDist}{Gaussian normal distribution}
\myNotationTableEntryMath{\ChiSqDist_{n}}{n-dimensional chi-square distribution}
%
\end{myNotationTable}


%%%%%%%%%%%%%%%%%%%%%%%%%%%%%%%%%%%%%%%%%%%%%%%%%%%%%%%%%%%%
\section*{Numbers and indexing}
%%%%%%%%%%%%%%%%%%%%%%%%%%%%%%%%%%%%%%%%%%%%%%%%%%%%%%%%%%%%
\begin{myNotationTable}[0.3\linewidth] %Breite der ersten Spalte anpassbar, sollte aber überall gleich sein
%
% Als zweispaltige Tabelle angelegt,
% Beim Makro myNotationTableEntryMath ist die erste Spalte bereits im Mathe-Modus
\myNotationTableEntryMath{\NatNum}{natural numbers}
\myNotationTableEntryMath{\NatNum_0}{natural numbers including zero (non-negative integers)}
\myNotationTableEntryMath{k, t}{discrete points in time}
\myNotationTableEntryMath{i, j, \ell, q}{indexing for objects, measurements and points}
%\myNotationTableEntryMath{m, n}{number of detections/measurements, number of tracked objects}
%
\end{myNotationTable}
%
%
%
%%%%%%%%%%%%%%%%%%%%%%%%%%%%%%%%%%%%%%%%%%%%%%%%%%%%%%%%%%%%%
%\section*{Indexing}
%%%%%%%%%%%%%%%%%%%%%%%%%%%%%%%%%%%%%%%%%%%%%%%%%%%%%%%%%%%%%
%\begin{myNotationTable}
%
%% Als zweispaltige Tabelle angelegt,
%% Beim Makro myNotationTableEntryMath ist die erste Spalte bereits im Mathe-Modus
%\myNotationTableEntryText{$k$, $t$}{discrete points in time}
%\myNotationTableEntryText{$i$, $j$, $q$}{indexing for objects, measurements and points}
%\myNotationTableEntryText{$m$, $n$}{number of detections/measurements, number of tracked objects}
%
%\end{myNotationTable}


%%%%%%%%%%%%%%%%%%%%%%%%%%%%%%%%%%%%%%%%%%%%%%%%%%%%%%%%%%%%
\section*{Geometry (coordinates, vehicle, {\newline}and camera modeling)}
%%%%%%%%%%%%%%%%%%%%%%%%%%%%%%%%%%%%%%%%%%%%%%%%%%%%%%%%%%%%
\begin{myNotationTable}[0.3\linewidth] %Breite der ersten Spalte anpassbar, sollte aber überall gleich sein
%
% Als zweispaltige Tabelle angelegt,
% Beim Makro myNotationTableEntryMath ist die erste Spalte bereits im Mathe-Modus
\myNotationTableEntryMath{x, y, z}{world coordinates}
\myNotationTableEntryMath{u, v}{image coordinates}
\myNotationTableEntryMath{b}{stereo base line}
\myNotationTableEntryMath{f}{focal length}
\myNotationTableEntryMath{\Delta{u}}{displacement in the image, disparity}
\myNotationTableEntryMath{d(\cdot)}{distortion function}
\myNotationTableEntryMath{\kappa_1, \kappa_2}{radial distortion parameters}
\myNotationTableEntryMath{\rho_1, \rho_2}{tangential distortion parameters}
%
\myNotationTableEntryMath{l, w, h}{length, width, height}
\myNotationTableEntryMath{r}{radius}
%
\myNotationTableEntryMath{A}{area}
\myNotationTableEntryMath{V}{volume}
%
\myNotationTableEntryMath{v}{velocity}
\myNotationTableEntryMath{a}{acceleration}
\myNotationTableEntryMath{\alpha}{steering angle}
%\myNotationTableEntryText{}{}
\myNotationTableEntryMath{\varphi}{orientation angle}
\myNotationTableEntryMath{\dot{\varphi}}{yaw rate}
%
\myNotationTableEntryMath{\pkt}{point in 2D and 3D space}
\myNotationTableEntryMath{\pHomogen}{point in homogeneous coordinates}
%
\myNotationTableEntryMath{\RotMat}{rotation matrix}
\myNotationTableEntryMath{\transVec}{translation vector}
%
\end{myNotationTable}


%%%%%%%%%%%%%%%%%%%%%%%%%%%%%%%%%%%%%%%%%%%%%%%%%%%%%%%%%%%%
\section*{Object state modeling and probabilities}
%%%%%%%%%%%%%%%%%%%%%%%%%%%%%%%%%%%%%%%%%%%%%%%%%%%%%%%%%%%%
\begin{myNotationTable}[0.3\linewidth] %Breite der ersten Spalte anpassbar, sollte aber überall gleich sein
%
% Als zweispaltige Tabelle angelegt,
% Beim Makro myNotationTableEntryMath ist die erste Spalte bereits im Mathe-Modus
\myNotationTableEntryMath{\statespace}{state space}
\myNotationTableEntryMath{\measurementspace}{measurement space}
%\myNotationTableEntryText{}{}
\myNotationTableEntryMath{\Fmat}{system matrix of the Kalman Filter}
\myNotationTableEntryMath{\Gmat}{control matrix of the Kalman Filter}
\myNotationTableEntryMath{\Hmat}{measurement matrix of the Kalman Filter}
\myNotationTableEntryMath{\KGk}{Kalman gain at time $k$}
%
\end{myNotationTable}
%
\end{showNotation}
%
%
%
%%%%%%%%%%%%%%%%%%%%%%%%%%%%%%%%%%%%%%%%
%%%%     automatisch erstelltes     %%%%
%%%%       Symbolverzeichnis        %%%%
%%%%%%%%%%%%%%%%%%%%%%%%%%%%%%%%%%%%%%%%
\begin{showListOfSymbols}%
\IfDefined{printglossary}{%
	\setlength{\LTleft}{0pt}% Einzug entfernen (Forderung des KSP-Verlages)
	%\setlength{\LTleft}{-5pt}% Bei Verwendung von style=long-booktabs
	%printsymbols[style=mylongglossstyle]
	%\printglossary[type=symbols,style=mylongglossstyle]
	%% Falls das Symbolverzeichnis eigene Breiten der Namen-Spalten haben soll:
	%\glssetwidest[0]{Haupteintrag}% <--- Hier längstes Wort unter den Glossar-Einträgen angeben
	%\glssetwidest[1]{Untereintraglänge}% <--- Hier längstes Wort unter den Glossar-Untereinträgen der 1 Ebene
	%\glssetwidest[2]{Untereintraglänge}% <--- Hier längstes Wort unter den Glossar-Untereinträgen der 2 Ebene
	\printglossary[type=symbols,style=alttree]
	%\printglossary[type=notation,style=longheader,sort=def] %,style=long-booktabs]
}%
\cleardoublepage
\end{showListOfSymbols}%
%
% Ab hier der Hauptteil:
\mainmatter
%